\documentclass[12pt,letterpaper]{article}

% ─── Paquetes ───────────────────────────────────────────────────────────────
\usepackage[utf8]{inputenc}
\usepackage[T1]{fontenc}
\usepackage[spanish]{babel}
\usepackage[margin=2.5cm]{geometry}
\usepackage{graphicx}
\usepackage{xcolor}
\usepackage{booktabs}
\usepackage{array}
\usepackage{tabularx}
\usepackage{enumitem}
\usepackage{listings}
\usepackage{fancyhdr}
\usepackage{titlesec}
\usepackage{hyperref}
\usepackage{parskip}
\usepackage{amssymb}
\usepackage{pifont}
\usepackage{tcolorbox}
\tcbuselibrary{skins, breakable}

% ─── Colores ────────────────────────────────────────────────────────────────
\definecolor{tecblue}{RGB}{0,70,127}
\definecolor{tecred}{RGB}{200,0,0}
\definecolor{codebg}{RGB}{28,35,52}
\definecolor{codebg2}{RGB}{245,247,250}
\definecolor{codeframe}{RGB}{100,120,160}
\definecolor{keywordcolor}{RGB}{79,140,255}
\definecolor{stringcolor}{RGB}{52,211,153}
\definecolor{commentcolor}{RGB}{138,144,153}
\definecolor{accentgreen}{RGB}{52,211,153}
\definecolor{accentyellow}{RGB}{251,191,36}
\definecolor{accentred}{RGB}{248,113,113}
\definecolor{accentpurple}{RGB}{167,139,250}

% ─── Encabezado / pie de página ─────────────────────────────────────────────
\pagestyle{fancy}
\fancyhf{}
\fancyhead[L]{\small CE-5507 · LinProd}
\fancyhead[R]{\small Guía de Implementación · Módulo 4}
\fancyfoot[C]{\thepage}
\renewcommand{\headrulewidth}{0.4pt}

% ─── Secciones ──────────────────────────────────────────────────────────────
\titleformat{\section}{\bfseries\large\color{tecblue}}{}{0em}{}[\titlerule]
\titleformat{\subsection}{\bfseries\normalsize\color{tecblue!70}}{}{0em}{}
\titleformat{\subsubsection}{\bfseries\normalsize}{}{0em}{}

% ─── Listings ───────────────────────────────────────────────────────────────
\lstset{
  language=Python,
  backgroundcolor=\color{codebg},
  basicstyle=\ttfamily\small\color{white},
  keywordstyle=\color{keywordcolor}\bfseries,
  stringstyle=\color{stringcolor},
  commentstyle=\color{commentcolor}\itshape,
  numberstyle=\tiny\color{commentcolor},
  numbers=left,
  numbersep=10pt,
  showstringspaces=false,
  breaklines=true,
  frame=single,
  framerule=0.4pt,
  rulecolor=\color{codeframe},
  tabsize=4,
  xleftmargin=15pt,
  xrightmargin=5pt,
  aboveskip=10pt,
  belowskip=10pt,
  literate=
    {á}{{\'a}}1 {é}{{\'e}}1 {í}{{\'i}}1 {ó}{{\'o}}1 {ú}{{\'u}}1
    {Á}{{\'A}}1 {É}{{\'E}}1 {Í}{{\'I}}1 {Ó}{{\'O}}1 {Ú}{{\'U}}1
    {ñ}{{\~n}}1 {Ñ}{{\~N}}1
    {←}{{$\leftarrow$}}1 {→}{{$\rightarrow$}}1
    {─}{{\textendash}}1 {—}{{\textemdash}}1,
}

% ─── Cajas especiales ───────────────────────────────────────────────────────
\tcbset{
  notebox/.style={
    enhanced, breakable,
    colback=tecblue!8,
    colframe=tecblue,
    fonttitle=\bfseries,
    attach boxed title to top left={yshift=-2mm, xshift=6mm},
    boxed title style={colback=tecblue, colframe=tecblue},
    title={#1},
  },
  warnbox/.style={
    enhanced,
    colback=accentyellow!10,
    colframe=accentyellow!80!black,
    fonttitle=\bfseries,
    attach boxed title to top left={yshift=-2mm, xshift=6mm},
    boxed title style={colback=accentyellow!80!black},
    title={#1},
  },
  tipbox/.style={
    enhanced,
    colback=accentgreen!8,
    colframe=accentgreen!70!black,
    fonttitle=\bfseries,
    attach boxed title to top left={yshift=-2mm, xshift=6mm},
    boxed title style={colback=accentgreen!70!black},
    title={#1},
  },
}

% ─── Checkmark ──────────────────────────────────────────────────────────────
\newcommand{\cmark}{\textcolor{accentgreen}{\ding{51}}}
\newcommand{\xmark}{\textcolor{accentred}{\ding{55}}}

% ─── Columna centrada ───────────────────────────────────────────────────────
\newcolumntype{C}[1]{>{\centering\arraybackslash}p{#1}}

% ═══════════════════════════════════════════════════════════════════════════
\begin{document}

% ─── Portada ────────────────────────────────────────────────────────────────
\begin{center}
  {\large\bfseries Proyecto 2 — LinProd}\\[0.5em]
  {\LARGE\bfseries\color{tecred} Guía de Implementación}\\[0.3em]
  {\LARGE\bfseries\color{tecblue} Módulo 4: Reportería PDF}\\[1em]
  \textit{Instituto Tecnológico de Costa Rica}\\
  \textit{Escuela de Ingeniería en Computadores}\\
  \textit{CE-5507 — Modelación Hardware Software Orientado a Objetos}\\
  \textit{I Semestre 2026}\\[1em]
  \textbf{Archivo a implementar:} \texttt{linprod/reports.py}
\end{center}

\vspace{0.5em}\hrule\vspace{1em}

% ─── 1. Contexto ────────────────────────────────────────────────────────────
\section{Contexto y rol del Módulo 4}

LinProd está dividido en cuatro módulos con una cadena de dependencia estricta
en un solo sentido:

\vspace{0.5em}
\begin{center}
\texttt{model.py} $\longleftarrow$ \texttt{simulator.py} $\longleftarrow$
\begin{tabular}{l}
  \texttt{app.py}\\
  \textbf{\texttt{reports.py} \textcolor{tecred}{$\longleftarrow$ tú aquí}}
\end{tabular}
\end{center}
\vspace{0.5em}

\textbf{Tu trabajo consiste únicamente en implementar \texttt{linprod/reports.py}}.
No debes modificar ningún otro archivo. El módulo recibe los datos del simulador
mediante la función \texttt{Simulator.get\_stats()} y genera un reporte en PDF
usando la librería \texttt{reportlab}.

\begin{tcolorbox}[notebox={Regla de oro}]
  \texttt{reports.py} \textbf{solo importa de} \texttt{linprod.simulator}
  (y por transitividad de \texttt{linprod.model}). No importa nada de
  \texttt{app.py} ni de \texttt{tkinter}.
\end{tcolorbox}

% ─── 2. Instalación ─────────────────────────────────────────────────────────
\section{Instalación de reportlab}

\texttt{reportlab} es la única dependencia no-stdlib del proyecto.
Para instalarla ejecuta en la terminal:

\begin{lstlisting}[language=bash, numbers=none]
pip install reportlab
# o si el sistema usa pip3:
pip3 install reportlab
\end{lstlisting}

Para verificar que quedó instalada correctamente:

\begin{lstlisting}[language=bash, numbers=none]
python3 -c "import reportlab; print(reportlab.Version)"
# Salida esperada: algo como 4.x.x
\end{lstlisting}

% ─── 3. Datos de entrada ─────────────────────────────────────────────────────
\section{Datos de entrada: \texttt{Simulator.get\_stats()}}

La función \texttt{get\_stats()} del simulador es la \textbf{única fuente de datos}
para el reporte. Retorna un diccionario con la siguiente estructura:

\subsection{Claves del diccionario principal}

\begin{center}
\begin{tabularx}{\textwidth}{|l|l|X|}
\hline
\textbf{Clave} & \textbf{Tipo} & \textbf{Descripción} \\
\hline
\texttt{total\_cycles\_run} & \texttt{int} & Ciclos totales ejecutados hasta terminar \\
\hline
\texttt{products\_total} & \texttt{int} & Total de productos que entraron a la línea \\
\hline
\texttt{products\_completed} & \texttt{int} & Productos que completaron toda la línea \\
\hline
\texttt{first\_completion\_cycle} & \texttt{int|None} & Ciclo en que salió el primer producto \\
\hline
\texttt{last\_completion\_cycle} & \texttt{int|None} & Ciclo en que salió el último producto \\
\hline
\texttt{avg\_completion\_time} & \texttt{float} & Tiempo promedio de paso (ciclos) \\
\hline
\texttt{total\_processing\_time} & \texttt{int} & Suma de tiempos de todos los productos \\
\hline
\texttt{bottleneck\_task} & \texttt{str|None} & Nombre de la tarea cuello de botella \\
\hline
\texttt{bottleneck\_process} & \texttt{str|None} & Proceso al que pertenece el cuello de botella \\
\hline
\texttt{avg\_wait\_to\_start} & \texttt{float} & Espera promedio global antes de iniciar tarea \\
\hline
\texttt{worst\_wait\_task} & \texttt{str|None} & Tarea con mayor espera promedio \\
\hline
\texttt{worst\_wait\_process} & \texttt{str|None} & Proceso de la tarea con mayor espera \\
\hline
\texttt{worst\_wait\_value} & \texttt{float} & Valor de esa mayor espera (ciclos) \\
\hline
\texttt{task\_stats} & \texttt{list[dict]} & Estadísticas por tarea (ver \S3.2) \\
\hline
\texttt{product\_stats} & \texttt{list[dict]} & Estadísticas por producto (ver \S3.3) \\
\hline
\end{tabularx}
\end{center}

\subsection{Estructura de cada elemento en \texttt{task\_stats}}

Cada elemento de la lista \texttt{task\_stats} es un diccionario con:

\begin{center}
\begin{tabularx}{\textwidth}{|l|l|X|}
\hline
\textbf{Clave} & \textbf{Tipo} & \textbf{Descripción} \\
\hline
\texttt{task\_name} & \texttt{str} & Nombre de la tarea \\
\hline
\texttt{process\_name} & \texttt{str} & Nombre del proceso al que pertenece \\
\hline
\texttt{process\_time} & \texttt{int} & Tiempo de procesamiento configurado (ciclos) \\
\hline
\texttt{products\_processed} & \texttt{int} & Cuántos productos procesó esta tarea \\
\hline
\texttt{avg\_wait\_cycles} & \texttt{float} & Espera promedio de productos en la cola \\
\hline
\texttt{busy} & \texttt{bool} & Si la tarea estaba ocupada al terminar \\
\hline
\end{tabularx}
\end{center}

\subsection{Estructura de cada elemento en \texttt{product\_stats}}

\begin{center}
\begin{tabularx}{\textwidth}{|l|l|X|}
\hline
\textbf{Clave} & \textbf{Tipo} & \textbf{Descripción} \\
\hline
\texttt{product\_id} & \texttt{int} & Identificador del producto (1, 2, 3...) \\
\hline
\texttt{entry\_time} & \texttt{int} & Ciclo en que entró a la línea \\
\hline
\texttt{exit\_time} & \texttt{int|None} & Ciclo en que salió (\texttt{None} si no terminó) \\
\hline
\texttt{is\_completed} & \texttt{bool} & Si completó toda la línea \\
\hline
\texttt{total\_time} & \texttt{int|None} & Ciclos totales de entrada a salida \\
\hline
\texttt{task\_history} & \texttt{list[dict]} & Historial de paso por cada tarea \\
\hline
\end{tabularx}
\end{center}

El historial de tareas (\texttt{task\_history}) de cada producto es una lista de
diccionarios, uno por cada tarea que procesó. Cada entrada tiene:
\texttt{task\_name}, \texttt{process\_name}, \texttt{start\_cycle},
\texttt{end\_cycle}, \texttt{wait\_cycles}.

% ─── 4. Qué debe contener el reporte ────────────────────────────────────────
\section{Contenido mínimo del reporte (según el enunciado)}

El enunciado del proyecto exige que el reporte presente los siguientes datos:

\begin{enumerate}[leftmargin=2em]
  \item[\cmark] Tiempo en que el \textbf{primer producto} logró terminar con
    éxito toda la línea (\texttt{first\_completion\_cycle}).
  \item[\cmark] Tiempo en que el \textbf{último producto} logró terminar con
    éxito toda la línea (\texttt{last\_completion\_cycle}).
  \item[\cmark] \textbf{Tiempo promedio} de paso por toda la línea
    (\texttt{avg\_completion\_time}).
  \item[\cmark] Proceso que generó el mayor \textbf{congestionamiento}
    (\textbf{cuello de botella}: \texttt{bottleneck\_process} /
    \texttt{bottleneck\_task}).
  \item[\cmark] \textbf{Tiempo promedio de espera} de productos para iniciar
    una tarea (\texttt{avg\_wait\_to\_start}).
  \item[\cmark] Proceso y tarea con \textbf{mayor tiempo de espera}
    (\texttt{worst\_wait\_process} / \texttt{worst\_wait\_task}).
  \item[\cmark] \textbf{Tiempo total} del procesamiento de todos los productos
    (\texttt{total\_processing\_time}).
  \item[\cmark] Cualquier otra estadística que el grupo considere necesaria.
\end{enumerate}

Se recomienda adicionalmente incluir la \textbf{tabla de detalle por tarea}
(\texttt{task\_stats}) y la \textbf{tabla de detalle por producto}
(\texttt{product\_stats}).

% ─── 5. Estructura del archivo ──────────────────────────────────────────────
\section{Estructura recomendada de \texttt{reports.py}}

\begin{lstlisting}
"""
reports.py — Módulo 4: Generación de reportes PDF
Proyecto LinProd | CE-5507 | I Semestre 2026
"""

from __future__ import annotations
from typing import Optional

from reportlab.lib.pagesizes import letter
from reportlab.lib.styles import getSampleStyleSheet, ParagraphStyle
from reportlab.lib.units import cm
from reportlab.lib import colors
from reportlab.platypus import (
    SimpleDocTemplate, Paragraph, Spacer,
    Table, TableStyle, HRFlowable,
)
from reportlab.lib.enums import TA_CENTER, TA_LEFT


def generate_report(stats: dict, output_path: str) -> None:
    """
    Genera el reporte PDF de la simulación.

    Parámetros
    ----------
    stats : dict
        Diccionario devuelto por Simulator.get_stats().
    output_path : str
        Ruta completa del archivo PDF a generar.
        Ejemplo: "/home/usuario/reporte_linprod.pdf"
    """
    doc = SimpleDocTemplate(
        output_path,
        pagesize=letter,
        leftMargin=2.5 * cm,
        rightMargin=2.5 * cm,
        topMargin=2.5 * cm,
        bottomMargin=2.5 * cm,
    )

    story = []           # lista de elementos (flowables)
    styles = _build_styles()

    # 1. Portada / encabezado
    _add_header(story, styles, stats)

    # 2. Resumen general
    _add_summary(story, styles, stats)

    # 3. Cuello de botella y esperas
    _add_bottleneck(story, styles, stats)

    # 4. Tabla detalle por tarea
    _add_task_table(story, styles, stats)

    # 5. Tabla detalle por producto
    _add_product_table(story, styles, stats)

    doc.build(story)


# ─── Helpers internos ────────────────────────────────────────────────────────

def _build_styles() -> dict:
    """Retorna un diccionario con los estilos de párrafo."""
    base = getSampleStyleSheet()
    return {
        "title":    base["Title"],
        "heading":  base["Heading2"],
        "body":     base["Normal"],
        "centered": ParagraphStyle(
            "centered", parent=base["Normal"], alignment=TA_CENTER
        ),
    }
\end{lstlisting}

\begin{tcolorbox}[notebox={¿Por qué este diseño?}]
  Dividir la generación en funciones privadas \texttt{\_add\_*} hace que el
  código sea fácil de depurar: si una sección da error, sabes exactamente
  dónde buscar.
\end{tcolorbox}

% ─── 6. Implementar cada sección ────────────────────────────────────────────
\section{Implementar cada sección del reporte}

\subsection{6.1 Encabezado y datos generales}

\begin{lstlisting}
def _add_header(story, styles, stats):
    story.append(Paragraph("LinProd — Reporte de Simulación", styles["title"]))
    story.append(HRFlowable(width="100%", thickness=1,
                             color=colors.HexColor("#4f8cff")))
    story.append(Spacer(1, 0.4 * cm))

    resumen = [
        ["Ciclos totales ejecutados:",
         str(stats["total_cycles_run"])],
        ["Productos en la línea:",
         str(stats["products_total"])],
        ["Productos completados:",
         str(stats["products_completed"])],
    ]
    t = Table(resumen, colWidths=[9 * cm, 7 * cm])
    t.setStyle(TableStyle([
        ("FONT",      (0, 0), (-1, -1), "Helvetica"),
        ("FONTSIZE",  (0, 0), (-1, -1), 11),
        ("FONTNAME",  (0, 0), (0, -1),  "Helvetica-Bold"),
        ("TEXTCOLOR", (0, 0), (0, -1),  colors.HexColor("#004680")),
        ("BOTTOMPADDING", (0, 0), (-1, -1), 5),
    ]))
    story.append(t)
    story.append(Spacer(1, 0.6 * cm))
\end{lstlisting}

\subsection{6.2 Resumen de tiempos}

\begin{lstlisting}
def _add_summary(story, styles, stats):
    story.append(Paragraph("Resumen de Tiempos", styles["heading"]))
    story.append(Spacer(1, 0.2 * cm))

    def fmt(v):
        return "—" if v is None else str(v)

    data = [
        ["Métrica", "Valor"],
        ["Primer producto completado (ciclo)",
         fmt(stats["first_completion_cycle"])],
        ["Último producto completado (ciclo)",
         fmt(stats["last_completion_cycle"])],
        ["Tiempo promedio de paso",
         f"{stats['avg_completion_time']:.2f} ciclos"],
        ["Tiempo total de procesamiento",
         f"{stats['total_processing_time']} ciclos"],
        ["Espera promedio global",
         f"{stats['avg_wait_to_start']:.2f} ciclos"],
    ]

    t = Table(data, colWidths=[10 * cm, 6 * cm])
    t.setStyle(TableStyle([
        # Encabezado
        ("BACKGROUND",   (0, 0), (-1, 0), colors.HexColor("#004680")),
        ("TEXTCOLOR",    (0, 0), (-1, 0), colors.white),
        ("FONTNAME",     (0, 0), (-1, 0), "Helvetica-Bold"),
        ("FONTSIZE",     (0, 0), (-1, 0), 11),
        # Cuerpo
        ("FONTNAME",     (0, 1), (-1, -1), "Helvetica"),
        ("FONTSIZE",     (0, 1), (-1, -1), 10),
        ("ROWBACKGROUNDS", (0, 1), (-1, -1),
         [colors.HexColor("#f0f4ff"), colors.white]),
        ("GRID", (0, 0), (-1, -1), 0.5, colors.HexColor("#cccccc")),
        ("BOTTOMPADDING", (0, 0), (-1, -1), 5),
        ("TOPPADDING",    (0, 0), (-1, -1), 5),
    ]))
    story.append(t)
    story.append(Spacer(1, 0.6 * cm))
\end{lstlisting}

\subsection{6.3 Cuello de botella}

\begin{lstlisting}
def _add_bottleneck(story, styles, stats):
    story.append(Paragraph("Cuello de Botella", styles["heading"]))
    story.append(Spacer(1, 0.2 * cm))

    bneck_proc = stats.get("bottleneck_process") or "—"
    bneck_task = stats.get("bottleneck_task") or "—"
    worst_val  = stats.get("worst_wait_value", 0.0)
    worst_task = stats.get("worst_wait_task") or "—"

    texto = (
        f"El cuello de botella detectado es la tarea <b>{bneck_task}</b> "
        f"del proceso <b>{bneck_proc}</b>, que registró la mayor "
        f"espera promedio en cola ({worst_val:.2f} ciclos por producto)."
    )
    story.append(Paragraph(texto, styles["body"]))
    story.append(Spacer(1, 0.6 * cm))
\end{lstlisting}

\subsection{6.4 Tabla de detalle por tarea}

\begin{lstlisting}
def _add_task_table(story, styles, stats):
    story.append(Paragraph("Detalle por Tarea", styles["heading"]))
    story.append(Spacer(1, 0.2 * cm))

    encabezado = [["Proceso", "Tarea", "t (ciclos)",
                   "Procesados", "Espera prom."]]
    filas = []
    for ts in stats["task_stats"]:
        filas.append([
            ts["process_name"],
            ts["task_name"],
            str(ts["process_time"]),
            str(ts["products_processed"]),
            f"{ts['avg_wait_cycles']:.2f}",
        ])

    data = encabezado + filas
    col_w = [4*cm, 4*cm, 2.5*cm, 2.5*cm, 3*cm]
    t = Table(data, colWidths=col_w)
    t.setStyle(TableStyle([
        ("BACKGROUND",   (0, 0), (-1, 0), colors.HexColor("#004680")),
        ("TEXTCOLOR",    (0, 0), (-1, 0), colors.white),
        ("FONTNAME",     (0, 0), (-1, 0), "Helvetica-Bold"),
        ("FONTNAME",     (0, 1), (-1, -1), "Helvetica"),
        ("FONTSIZE",     (0, 0), (-1, -1), 9),
        ("ROWBACKGROUNDS", (0, 1), (-1, -1),
         [colors.HexColor("#f0f4ff"), colors.white]),
        ("GRID", (0, 0), (-1, -1), 0.5, colors.HexColor("#cccccc")),
        ("ALIGN",    (2, 0), (-1, -1), "CENTER"),
        ("BOTTOMPADDING", (0, 0), (-1, -1), 4),
        ("TOPPADDING",    (0, 0), (-1, -1), 4),
    ]))
    story.append(t)
    story.append(Spacer(1, 0.6 * cm))
\end{lstlisting}

\subsection{6.5 Tabla de detalle por producto}

\begin{lstlisting}
def _add_product_table(story, styles, stats):
    story.append(Paragraph("Detalle por Producto", styles["heading"]))
    story.append(Spacer(1, 0.2 * cm))

    encabezado = [["Producto", "Entrada", "Salida",
                   "Total (ciclos)", "Estado"]]
    filas = []
    for p in stats["product_stats"]:
        filas.append([
            f"P{p['product_id']}",
            str(p["entry_time"]),
            str(p["exit_time"]) if p["exit_time"] is not None else "—",
            str(p["total_time"]) if p["total_time"] is not None else "—",
            "Completado" if p["is_completed"] else "En proceso",
        ])

    data = encabezado + filas
    col_w = [2.5*cm, 2.5*cm, 2.5*cm, 3.5*cm, 5*cm]
    t = Table(data, colWidths=col_w)
    t.setStyle(TableStyle([
        ("BACKGROUND",   (0, 0), (-1, 0), colors.HexColor("#004680")),
        ("TEXTCOLOR",    (0, 0), (-1, 0), colors.white),
        ("FONTNAME",     (0, 0), (-1, 0), "Helvetica-Bold"),
        ("FONTNAME",     (0, 1), (-1, -1), "Helvetica"),
        ("FONTSIZE",     (0, 0), (-1, -1), 9),
        ("ROWBACKGROUNDS", (0, 1), (-1, -1),
         [colors.HexColor("#f0f4ff"), colors.white]),
        ("GRID", (0, 0), (-1, -1), 0.5, colors.HexColor("#cccccc")),
        ("ALIGN",    (1, 0), (-1, -1), "CENTER"),
        ("BOTTOMPADDING", (0, 0), (-1, -1), 4),
        ("TOPPADDING",    (0, 0), (-1, -1), 4),
    ]))
    story.append(t)
\end{lstlisting}

% ─── 7. Cómo integrarlo en la GUI ───────────────────────────────────────────
\section{Integración con la GUI (\texttt{app.py})}

El botón \textbf{«= Estadísticas»} ya existe en la barra de controles.
Para agregar la generación de PDF hay que añadir un botón \textbf{«PDF»}
en esa misma barra y conectarlo al siguiente handler:

\begin{lstlisting}
# En app.py, importar al inicio:
from tkinter import filedialog
from linprod.reports import generate_report

# Handler del nuevo botón:
def _on_generate_pdf(self) -> None:
    if self._sim is None or not self._sim.is_started:
        messagebox.showinfo("Sin datos",
                            "Inicie y complete una simulación primero.")
        return
    stats = self._sim.get_stats()
    path = filedialog.asksaveasfilename(
        defaultextension=".pdf",
        filetypes=[("PDF", "*.pdf")],
        initialfile="reporte_linprod.pdf",
    )
    if not path:
        return   # usuario canceló
    try:
        generate_report(stats, path)
        messagebox.showinfo("PDF generado",
                            f"Reporte guardado en:\n{path}")
    except Exception as e:
        messagebox.showerror("Error al generar PDF", str(e))
\end{lstlisting}

Y en \texttt{\_build\_controls\_card}, después del botón de estadísticas:

\begin{lstlisting}
self._btn_pdf = ModernButton(btns, text="PDF  Generar PDF",
                             variant="secondary",
                             command=self._on_generate_pdf)
self._btn_pdf.pack(side="left", padx=6)
\end{lstlisting}

% ─── 8. Prueba sin GUI ──────────────────────────────────────────────────────
\section{Prueba del reporte sin necesidad de la GUI}

Para probar \texttt{reports.py} de forma independiente, sin abrir la interfaz
gráfica, puedes agregar un bloque \texttt{if \_\_name\_\_ == "\_\_main\_\_":}
al final del archivo:

\begin{lstlisting}
if __name__ == "__main__":
    from linprod.model import Process, ProductionLine, Task
    from linprod.simulator import Simulator

    # Línea de prueba mínima
    line = ProductionLine()

    p1 = Process("Corte",   is_initial=True)
    p1.add_task(Task("Corte láser", 2))
    p1.add_task(Task("Pulido",      2))

    p2 = Process("Ensamble")
    p2.add_task(Task("Armado",    3))
    p2.add_task(Task("Soldadura", 4))

    p3 = Process("Empaque", is_final=True)
    p3.add_task(Task("Empacar",    2))
    p3.add_task(Task("Etiquetado", 1))

    line.add_process(p1)
    line.add_process(p2)
    line.add_process(p3)

    sim = Simulator(line, n_products=5)
    sim.start()
    sim.run()    # corre todos los ciclos hasta terminar

    stats = sim.get_stats()
    generate_report(stats, "reporte_prueba.pdf")
    print("PDF generado: reporte_prueba.pdf")
\end{lstlisting}

Para ejecutarlo desde la raíz del repositorio:

\begin{lstlisting}[language=bash, numbers=none]
python3 -m linprod.reports
\end{lstlisting}

% ─── 9. Errores comunes ─────────────────────────────────────────────────────
\section{Errores comunes y cómo resolverlos}

\begin{center}
\begin{tabularx}{\textwidth}{|p{5.5cm}|X|}
\hline
\textbf{Error} & \textbf{Solución} \\
\hline
\texttt{ModuleNotFoundError: reportlab} &
Instalar con \texttt{pip3 install reportlab} \\
\hline
\texttt{PermissionError} al guardar el PDF &
El archivo ya está abierto en el visor. Ciérralo antes de generarlo de nuevo. \\
\hline
\texttt{KeyError: 'alguna\_clave'} &
Verificar que estás accediendo a la clave correcta del diccionario (revisar
\S3 de esta guía). \\
\hline
Las tildes se ven como \texttt{?} en el PDF &
Usar \texttt{pdfmetrics.registerFont} con una fuente TTF que soporte UTF-8,
o evitar tildes en los textos generados. \\
\hline
Tabla cortada al final de la página &
Usar \texttt{KeepTogether} para proteger tablas cortas, o \texttt{TableStyle}
con \texttt{splitByRow=1} para las largas. \\
\hline
\end{tabularx}
\end{center}

\begin{tcolorbox}[tipbox={Consejo sobre tildes}]
  La fuente por defecto de reportlab (\texttt{Helvetica}) no soporta tildes
  directamente. La solución más sencilla es usar la fuente del sistema
  \textbf{DejaVu} que sí soporta UTF-8:

\begin{lstlisting}[numbers=none, aboveskip=4pt, belowskip=0pt]
from reportlab.pdfbase import pdfmetrics
from reportlab.pdfbase.ttfonts import TTFont

pdfmetrics.registerFont(
    TTFont("DejaVu", "/usr/share/fonts/truetype/dejavu/DejaVuSans.ttf")
)
# Después usas "DejaVu" como nombre de fuente en TableStyle/ParagraphStyle
\end{lstlisting}
\end{tcolorbox}

% ─── 10. Lista de verificación ──────────────────────────────────────────────
\section{Lista de verificación antes de la entrega}

\begin{enumerate}[leftmargin=2em]
  \item[\cmark] \texttt{reports.py} importa únicamente de \texttt{linprod.simulator},
    \texttt{reportlab} y la stdlib. No toca \texttt{app.py}.
  \item[\cmark] La función pública es \texttt{generate\_report(stats, output\_path)}.
  \item[\cmark] El PDF incluye los 7 ítems obligatorios del enunciado (\S4).
  \item[\cmark] Las tablas tienen encabezado con fondo de color para legibilidad.
  \item[\cmark] Los valores \texttt{None} se muestran como ``—'' (no como la cadena
    \texttt{"None"}).
  \item[\cmark] El módulo puede ejecutarse de forma autónoma (\texttt{python3 -m linprod.reports})
    para pruebas sin GUI.
  \item[\cmark] Se probó con la línea demo de 5 procesos × 2 tareas × 5 productos.
  \item[\xmark] \textbf{No} modificar \texttt{model.py}, \texttt{simulator.py} ni
    \texttt{app.py} (salvo agregar el botón PDF indicado en \S7).
\end{enumerate}

% ─── Fin ────────────────────────────────────────────────────────────────────
\vspace{1em}
\hrule
\vspace{0.5em}
\begin{center}
  \small
  Proyecto 2 — LinProd \textbullet{}
  CE-5507 Modelación Hardware/Software OO \textbullet{}
  TEC, I Semestre 2026
\end{center}

\end{document}
